\documentclass[11pt]{article}

\usepackage[margin=1in]{geometry}
\usepackage{amsmath,amssymb}
\usepackage{array}
\usepackage{booktabs}
\usepackage{enumitem}
\usepackage{hyperref}
\usepackage{microtype}
\usepackage[T1]{fontenc}
\usepackage{lmodern}

\hypersetup{
    colorlinks=true,
    linkcolor=blue,
    urlcolor=blue
}

\title{\textit{cs-mail} Express Lanes --- Implementation Plan\\
\large Recipient-Granted, Communication for Chosen Organizations}
\author{C-SQD}
\date{Draft 0.4 --- September 2026}

\setlength{\emergencystretch}{1.5em}

\begin{document}
\maketitle

\begin{center}
\fbox{\parbox{0.92\textwidth}{\small\textbf{Nonnormative implementation plan.}
This document records settled design decisions and a build order for
\emph{express lanes}, the launch-feature form of the scoped-capability extension
that the \textit{cs-mail Protocol Specification} describes. It extends the
\textit{Rust Reference Architecture} as a target design for the revised request
model. Existing API names below identify adapter boundaries, not a claim that
the current workspace implements the revised semantics. The specification
controls admission and request outcomes; the C-SQD financial profile controls
the separate member program. The user-facing
explanation of the feature is the companion memo, \textit{Express Lanes for
Transactional Mail}.}}
\end{center}

\section{The Central Observation}

A lane grants bounded communication permission without a relationship-request
charge. Its admission path authenticates the sender, checks blocking and the
signed scope, validates durable content, consumes allowance, and records delivery.
It creates no refund or forfeiture. Accepted delivery and permitted follow-ups
also move no request value, although a follow-up must join an identified open
request under current recipient policy.

This distinction places lane evaluation at the durable admission boundary. It
does not imply that the request kernel runs only when money moves: relationship
transitions can change permission without a current charge. The broader business
revision changes request settlement; lanes preserve their own absence of financial
effects within that revised contract.

\section{Settled Design Decisions}

\begin{enumerate}[leftmargin=2em]
    \item \textbf{Edge subsystem, not kernel state.} A lane is evaluated at the
    admission edge, before the kernel is consulted---mirroring how accepted
    delivery already avoids the kernel. It uses a bounded, idempotent in-memory model for tests and a durable
    PostgreSQL coordinator for admitted delivery.
    \item \textbf{Domain as identity, for legacy senders.} A DMARC-authenticated
    sending domain resolves to a single synthetic sender protocol identity, so
    every authenticated address at that domain shares one lane and the per-pair
    aggregate is preserved.
    \item \textbf{One lane type, a subject enum.} A lane's subject is either a
    native opaque handle or a recipient-granted legacy domain scope. The handle
    identifies the pair; a separate recipient signature authorizes the complete
    grant. The same versioned record's subject can be swapped on adoption.
    \item \textbf{DMARC verification as evidence-as-data.} Verification happens at
    the edge and produces a \texttt{VerifiedDomain} value. A separately
    versioned binding joins that proof to the recipient-granted lane domain, so
    lane evaluation remains deterministic and testable without DNS.
    \item \textbf{Lifetime on the existing scheduler.} Lanes are time-bound on a
    sliding window anchored to last activity, bounded by the grant's hard
    \texttt{not\_after}, using the shared \texttt{ScheduleTask} infrastructure,
    \texttt{cs\_schedules}, and the \texttt{cs-mail-worker} lease.
    \item \textbf{Block and lane are mutually exclusive.} A directed pair is at
    most one of \{blocked, lane-authorized, neither\}. Blocking hard-revokes any
    lane; granting a lane to a blocked pair requires an explicit unblock;
    unblocking leaves the pair unaccepted without restoring a lane.
    \item \textbf{Issuance is recipient-initiated only.} Lanes are granted at
    address handover (checkout, signup, booking) or by replying to an
    organization first. The ``accept for receipts only'' first-contact path is
    removed.
    \item \textbf{Volume and declared scope are enforced.} Each lane carries a
    maximum message count per interval that admission consumes atomically, a
    versioned declared-purpose constraint, an optional origin constraint, and the
    declaration authorities it accepts. Admission compares authenticated
    declarations exactly. This proves what a sender or gateway declared; providers
    still cannot inspect encrypted content or certify that the declaration is
    semantically truthful.
    \item \textbf{Lane events do not settle requests.} Grant, renewal, lapse, and
    revocation change lane authority only. A separate pending request retains its
    deadline and financial outcome rules.
\end{enumerate}

\section{The Admission Decision}

Express lanes add one branch to the answer of a question the system already asks:
which authority permits this inbound message? The ordered decision, evaluated for a
directed \texttt{(sender, recipient)} pair at a single canonical boundary, is:

\begin{enumerate}[leftmargin=2em]
    \item if the pair is \texttt{Blocked}, refuse (\texttt{ContactBlocked});
    \item else if the relationship is \texttt{Accepted}, deliver uncharged through
    the existing accepted-sender path;
    \item else if a valid, in-window lane authorizes the pair and the authenticated
    message declarations match its constraints, deliver uncharged through the lane
    path;
    \item else, if the message identifies an open request, admit a follow-up only
    if current recipient policy permits it; refusal cannot create another charge;
    \item else, if a preparing or open request already exists, return its scoped
    state without creating a new request;
    \item otherwise offer an eligible new-request flow, with explicit sender
    authorization before any charge.
\end{enumerate}

The second and third branches share one durable uncharged admission path. A lane
admission creates no request charge, moves no ledger value, opens no
relationship solicitation episode, and advances no repeated-attempt state---the
same properties required of capability-authorized communication. Tests must
establish these effects at the storage boundary; avoiding a particular function
call is not by itself proof of the invariant.

\begin{verbatim}
enum AdmissionOutcome {
    Refused,          // blocked pair
    AcceptedPermission, // existing accepted-sender path
    LanePermission,     // new: valid express lane
    RequestFollowup, // open request and current recipient allowance
    ExistingRequest, // resume preparation or report follow-up refusal
    NewRequestOffer, // eligibility checked; explicit charge authorization needed
}
\end{verbatim}

For native mail, the sender signs the complete uncharged admission request,
including content reference, declared purpose, origin declaration, optional
payload-schema identifier, message validity, capability reference, evidence,
idempotency key, protocol version, deployment domain, and intended provider. The
same declarations are cryptographically bound to the encrypted content. A legacy
SMTP sender cannot supply a native protocol signature: the trusted ingress service
constructs and signs the request only after its configured \texttt{DmarcVerifier}
returns passing evidence, identifies itself as the declaration authority, and uses
legacy or unspecified origin unless it has separate authenticated evidence.

A scope mismatch returns \texttt{CapabilityScopeMismatch}, consumes no rate
allowance, and moves no value. The message may use another path only under its independently valid authority.
A pending request cannot be bypassed by buying another charge. Admission after
message validity returns \texttt{MessageValidityClosed}: it does not settle an
open request merely because one follow-up message is stale. A preparing initial
request that cannot be admitted instead follows the full cancellation/refund rule.

\section{The Lane Object and Its Store}

A lane is a compact recipient-granted record. It carries no money and references
no message content.

\begin{verbatim}
struct Lane {
    subject: LaneSubject,        // native identity or verified domain
    recipient: ProtocolIdentity,
    purpose: PurposeConstraint,  // matches authenticated declaration
    origin: Option<OriginConstraint>,
    declaration_authority: DeclarationAuthorityConstraint,
    rate_limit: RateLimit,        // enforced atomically at admission
    window: Duration,            // sliding lifetime
    not_after: CanonicalTime,     // signed hard validity bound
    last_activity: CanonicalTime,
    mode: LaneMode,              // Expiring | Persistent
    version: Version,
}

enum LaneSubject {
    Native(ScopedHandle),        // opaque identifier from cs-mail-privacy
    LegacyDomain(DomainIdentity), // recipient-granted domain scope
}
\end{verbatim}

The stored lane is created only from a \texttt{SignedLaneGrant}. The recipient
signature covers the subject, directed protocol pair, purpose and optional origin
constraints, accepted declaration authority, rate limit, sliding lifetime, hard
validity interval, mode, protocol version, deployment, intended provider,
operational key, and grant version. \texttt{ScopedHandle} is an opaque identifier,
not the signature or the authority by itself.

\texttt{DeclaredPurpose}, \texttt{OriginDeclaration}, declaration-authority, and
schema identifier are constrained, versioned domain values rather than free-form
strings. The base purpose vocabulary includes \texttt{Unspecified}; future
purposes and schemas use namespaced identifiers. Unknown noncritical schemas do
not prevent lane delivery, while an unsupported critical schema is refused before
allowance is consumed.

The lane store is a new, compartmentalized access domain: its own database table,
credentials, encryption keys, and access role, per the specification's
relationship-graph privacy requirements. It is not a column on the relationship
aggregate and not part of \texttt{protocol\_state}. A ledger or telemetry
compromise must not yield the list of organizations a recipient trusts. Grant,
renew, and revoke operations are idempotent record transitions with no ledger
coupling, exactly like the existing adapter books.

The storage coordinator enforces one invariant under the relationship aggregate
lock: for any directed pair, at most one of \{blocked, lane-authorized\} holds at
a time.

\section{Legacy Senders: Domain as Identity}

Most transactional senders will not be on \textit{cs-mail} at launch, so the
common case is legacy email authenticated by DKIM and DMARC. The inbound edge
verifies DMARC under RFC~9989 and produces a \texttt{VerifiedDomain}; lane evaluation
consumes that value and never performs DNS or signature work itself. This keeps
lane logic pure and unit-testable and confines the messy legacy-world
cryptography to one edge boundary.

{\sloppy
A passing result keeps the RFC 5322 author domain, independent DKIM and SPF pass
facts (each with its applicable strict or relaxed alignment mode), and canonical
evaluation time as evidence. It does not claim that the verifier exposed a DMARC
policy-record or organizational domain when the underlying library did not. Lane
identity is instead the canonical IDNA-ASCII domain scope deliberately selected
by the recipient's grant; it is never whichever DKIM or SPF identifier happened
to pass. The authenticated author must be that domain or a subdomain of it. That
lane domain resolves through a versioned, domain-separated mapping to
a \emph{single} synthetic sender protocol identity.
Consequently \texttt{statements@bank.com} and \texttt{alerts@bank.com} are one
sender to the system, one lane covers both, and admission still keys to a single
\texttt{(sender, recipient)} aggregate. A message that fails DMARC alignment to a
lane's domain is simply not on the lane, regardless of its display name or
local-part. DMARC resists unauthorized use of the expected domain; the separate
expected-domain match rejects look-alike domains.\par}

DMARC says nothing about message purpose, automation mode, or human authorship.
The gateway therefore marks the envelope as gateway-declared and uses
\texttt{LegacyOrUnspecified} origin unless another authenticated mechanism
supports a narrower statement. A recipient grant explicitly decides whether a
gateway-originated purpose declaration is acceptable. The user interface must not
present that choice as content verification; recipient reports, revocation, and
blocking remain the response to dishonest use of the declared scope.

The earlier adapter design identifies a dedicated \texttt{cs-mail-smtp-gateway} library. It specifies
Stalwart Labs' \texttt{mail-auth} 0.12.1 with default features disabled and the
minimal Ring/Hickory DNS feature set. The gateway accepts the remote IP,
EHLO identity, MAIL FROM (including the null reverse-path case), receiver hostname,
raw message, and canonical receipt time. It bounds message size, header size and
count, DKIM-signature count, and receipt/system-clock skew before performing
DNS-backed verification. Temporary DNS failures remain distinguishable from
authentication failure, malformed input, and resource exhaustion for SMTP policy.

\section{Native Senders: Pairwise Handles}

For a sender already on \textit{cs-mail}, the lane binds to the pairwise-scoped
\texttt{ScopedHandle} already provided by \texttt{cs-mail-privacy}. The recipient
signs the full lane grant containing that handle. The provider verifies the grant at admission
without learning which organization it authorizes, so the native lane is fully
private where the legacy lane must expose a domain. When a legacy sender adopts
\textit{cs-mail}, its lane's subject is swapped from \texttt{LegacyDomain} to
\texttt{Native}, upgrading the lane in place and preserving its history.

Native identity also enables finer control that domain granularity cannot offer.
Because a legacy domain is a single identity, block and lane are mutually
exclusive at domain granularity---one cannot block a bank's marketing while
keeping a lane for its fraud alerts. A native organization that uses distinct
protocol identities for distinct purposes can be blocked on one and lane-granted
on another.

\section{Lifetime and Renewal}

A lane's clock is its last authenticated message, not its issue time. Each
admitted lane message bumps \texttt{last\_activity} and reschedules the lane's
horizon; a lane in active use therefore never lapses on a timer, while a dormant
one drifts toward expiry.

Two modes cover the range the memo describes:

\begin{itemize}[leftmargin=2em]
    \item \texttt{Expiring} (default): at the horizon the recipient receives a
    gentle re-confirm prompt rather than a silent deletion. The window is set per
    sender.
    \item \texttt{Persistent}: the same renew-on-activity behavior, but the horizon
    keeps the lane active and emits a review reminder instead of lapsing it. The
    signed \texttt{not\_after} remains a hard stop. This is the mode for a
    bank's fraud alerts or a hospital's results, where a silently lapsed lane
    would be costly.
\end{itemize}

Both modes reuse the leased \texttt{cs\_schedules} plus \texttt{cs-mail-worker}
machinery: a scheduled task fires at \texttt{last\_activity + window} and either
lapses the lane or emits a review effect through the outbox. The generic
\texttt{ScheduleTask} and \texttt{ScheduleChange} types live in
\texttt{cs-mail-primitives}; no second lease substrate is required.

\section{Block and Lane Interaction}

Block and lane are opposite ends of one recipient control, so the model treats
them as mutually exclusive rather than as independent flags. The two directions
are deliberately asymmetric:

\begin{itemize}[leftmargin=2em]
    \item \textbf{Blocking hard-revokes the lane, unconditionally.} Block is the
    most restrictive state and a deliberate protective action, so it tears down
    any lane in the same commit that records the block. After a later unblock the
    organization does \emph{not} regain its lane; it returns to an unaccepted sender,
    subject to preserved request history and any other independently valid authority.
    \item \textbf{Granting a lane never silently unblocks.} Because issuance is
    one-tap or automatic, an issuance flow must not silently defeat a block the
    recipient deliberately set. A grant targeting a blocked pair surfaces the
    block and requires an explicit unblock-and-allow confirmation.
\end{itemize}

This mutual-exclusion invariant also resolves the coordination question raised by
placing lanes outside the kernel. Block state lives in the relationship aggregate;
the lane lives in the edge store. Rather than reconciling two independent truths,
blocking is one state transition that revokes the lane in the same commit, and
lane admission reads the pair's block state transactionally---by taking the
relationship aggregate's row lock, or by consulting a minimal transactional block
flag---so a block that commits first is always observed. Domain-as-identity is
what makes ``the pair'' well defined: blocking \texttt{bank.com} and holding a
lane for \texttt{bank.com} concern the same directed identity pair.

The reference implementation chooses the aggregate row lock. A future split-store
deployment must preserve the same linearization contract through a narrowly
authorized transactional coordinator.

\section{Issuance}

A lane is granted only by a deliberate recipient act:

\begin{itemize}[leftmargin=2em]
    \item \textbf{At address handover} (checkout, signup, booking). The client
    records the web origin it is on as the expected sending domain, then completes
    the grant on the first DMARC-authenticated message from that domain, showing a
    one-tap confirmation that names the \emph{authenticated} domain rather than
    the display name. Seeding the expectation from the origin turns ``which
    sender?'' into a match and lets the recipient catch a look-alike domain at
    grant time.
    \item \textbf{On replying first.} Writing to an organization auto-opens a
    return lane, since the recipient has plainly invited a reply.
\end{itemize}

The first-contact ``accept for receipts only'' path is removed. A cold sender submits
one eligible relationship request; it can be accepted as a full standing contact, rejected,
or left to expire, but it is not granted a lane from that interaction. The
consequence is minor and recoverable: an organization one would have preferred to
keep on a narrow lane must instead be reached through handover or a reply, or
accepted fully. Issuance neither charges the sender nor settles that request.

\section{Changes by Crate}

\begin{itemize}[leftmargin=2em]
    \item \textbf{\texttt{cs-mail-capabilities}} holds the lane ``book'': the
    lane record, subject enum, purpose and origin constraints, declaration-
    authority rules, and idempotent grant/renew/revoke transitions, with no ledger
    coupling and the block/lane mutual-exclusion invariant. Its admission check
    performs exact versioned declaration matching before consuming allowance.
    \item \textbf{\texttt{cs-mail-adapters}} defines the async, library-neutral
    SMTP authentication request, typed failures, and versioned evidence model,
    alongside the existing outbound downgrade route.
    \item \textbf{\texttt{cs-mail-smtp-gateway}} implements that boundary with
    pinned Stalwart \texttt{mail-auth}, minimal crypto/DNS features, resource
    limits, and receipt-time skew checks.
    \item \textbf{\texttt{cs-mail-storage-postgres}} gains a migration for the
    compartmentalized lane table under its own access role and keys, plus the
    transactional block/lane coordination.\par
    \item \textbf{\texttt{cs-mail-primitives}, \texttt{cs-mail-worker}, and
    \texttt{cs\_schedules}} supply shared task types and the existing lease for
    lane horizon tasks (lapse or review); no new scheduling substrate.
    \item \textbf{\texttt{cs-mail-privacy}} supplies \texttt{ScopedHandle} as the
    native pair identifier; \texttt{cs-mail-capabilities} supplies the separately
    signed authority.
    \item \textbf{\texttt{cs-mail-primitives}} defines constrained purpose,
    origin, declaration-authority, message-validity, schema, and extension-
    criticality values without acquiring workflow or policy behavior.
    \item \textbf{\texttt{cs-mail-content} and \texttt{cs-mail-wire}} bind the
    declarations, capability reference, and validity time to encrypted content and
    provide canonical, versioned representations. Unknown noncritical schemas
    remain opaque; unsupported critical schemas fail before admission.
    \item \textbf{Service and capability layers} add authenticated ingress and
    domain-separated, versioned signing representations for grants and uncharged
    admission. Provider admission policy consumes envelope-visible facts only;
    recipient presentation policy remains in the client and has no settlement
    authority.
    \item \textbf{\texttt{cs-mail-protocol::transition}} (the settlement kernel)
    owns request outcomes. Lane operations do not invent new request settlement,
    while shared admission coordination observes the revised request contract.
\end{itemize}

\section{Build Order}

Each lane stage is independently testable after the revised request baseline
is established. References to earlier implementation status are not conformance
evidence for this revision.

\begin{enumerate}[leftmargin=2em]
    \item \textbf{Native lanes, end to end.} The lane book, the admission branch,
    grant and revoke, and the block/lane mutual-exclusion invariant with its
    transactional coordination established the permission layer and its hardest
    correctness property first.
    \item \textbf{Lifetime and renewal.} Sliding window, the two modes, and the
    scheduled lapse/re-confirm tasks on the existing worker lease.
    \item \textbf{Legacy senders.} The inbound DMARC verifier,
    domain-to-identity resolution, and the legacy subject, with independent SPF
    and DKIM evidence and expected-domain mismatch enforcement.
    \item \textbf{Issuance flows.} Record the expected domain at address handover, confirm the first
    authenticated message, and authorize a return lane when the recipient writes first.
    \item \textbf{Authenticated declaration enforcement.} Bind constrained purpose and origin declarations to native
    content and uncharged admission signing bytes, add explicit gateway authority,
    and reject mismatches without consuming lane allowance.
\end{enumerate}

\section{Open Questions and Risks}

\begin{itemize}[leftmargin=2em]
    \item \textbf{Split-store linearization.} The reference implementation uses
    the relationship aggregate row lock. Any future split-store deployment must
    preserve that ordering contract with a narrowly authorized coordinator.
    \item \textbf{Renewal signal.} An authenticated, admitted lane message renews
    the sliding horizon using monotonic canonical time. Decide separately whether
    recipient reads or replies should count.
    \item \textbf{DMARC operations.} The reference selects pinned Stalwart
    \texttt{mail-auth} and keeps DNS out of lane logic. A deployment must still
    operate resolver policy, map typed failures to SMTP replies, monitor latency,
    and keep its system clock within the configured receipt-time skew. Senders
    without a DMARC pass receive no legacy lane; a weaker grant would be a
    separately named and reviewed feature.
    \item \textbf{Per-subject retention.} Native lane subjects are opaque; legacy
    subjects expose a domain. Retention and audit rules differ per subject and
    must be encoded in the store's retention classes.
    \item \textbf{Issuance as a surface.} Rate-limit issuance confirmations and
    prevent the lane store or its prompts from becoming an enumeration oracle.
    \item \textbf{Declaration honesty and metadata.} A matched declaration is not
    proof that ciphertext has the claimed meaning. Keep the vocabulary coarse,
    retain it only as long as required for admission and audit, expose no recipient
    policy beyond actionable constraints, and make report, revoke, and block easy.
\end{itemize}

\section{Request and Member-Program Integration}

Add conformance stories for a lane message while a relationship request is open,
lane revocation followed by an allowed or refused request follow-up, and a block
racing both paths. Granting a lane does not accept the relationship or refund its
request. Lane expiry cannot extend the request deadline, and an unsupported
declaration cannot silently start a charged request.

An admitted lane message may renew the lane's activity horizon. It must not count
automatically as intentional member activity, contribute forfeitures, or affect a
quarterly allocation. These programs may reuse scheduling and scoped identifiers,
but their authorities and financial effects remain separate.

The companion memo explains this behavior to users. The protocol specification
defines admission and request invariants; the deployment profile supplies the
member program. Update examples against those versions rather than relying on
an earlier assertion that all settlement semantics are unchanged.

\end{document}
